\chapter[1926 --- Fibiger]{1926: Johannes Fibiger}
\label{ch:1926_johannesfibiger}

\section*{Main Contribution}

\textit{Discovered that a parasitic worm (Spiroptera) could cause cancer in animals, providing the first experimental evidence that parasites could induce malignancy.}

\section{Official Citation}

for his discovery of the Spiroptera carcinoma

\section{Background and Historical Context}

Johannes Andreas Grib Fibiger was born on April 15, 1867, in Silkeborg, Denmark\index{Fibiger, Johannes|textbf}. He studied medicine at the University of Copenhagen, where he developed a strong interest in pathology and bacteriology---disciplines that were then undergoing rapid development in the wake of the germ theory of disease\index{Pathology|textbf}. After receiving his medical degree in 1890, Fibiger worked in the pathological laboratory of Professor Christian Gram (of Gram stain fame) and later spent time in Robert Koch's laboratory in Berlin, where he acquired advanced techniques in bacteriology and experimental pathology\index{Christian Gram|textbf}.

Upon returning to Denmark, Fibiger was appointed professor of pathological anatomy at the University of Copenhagen in 1900, a position he held for the remainder of his career\index{University of Copenhagen|textbf}. He built a thriving research department that attracted students and collaborators from across Europe, and he established a reputation as one of the leading experimental pathologists of his generation. His research focused on the relationship between irritation, inflammation, and cancer---a topic that was of intense interest in the early twentieth century, as the causes of cancer were almost entirely unknown.

Fibiger's Nobel Prize was awarded for what was, at the time, considered the first experimental demonstration that a parasite could cause cancer---a finding of enormous theoretical significance\index{Carcinogenesis|textbf}. However, the story of Fibiger's discovery, and the subsequent reassessment of his findings, is one of the most fascinating and instructive episodes in the history of cancer research.

The late nineteenth and early twentieth centuries were a period of intense interest in the etiology of cancer\index{Cancer research|textbf}. Several theories competed to explain the origin of malignant tumors. The parasitic theory of cancer, which had been advocated by several prominent pathologists, held that cancer was caused by infectious agents---parasites, bacteria, or viruses---that invaded tissues and triggered malignant transformation. This theory was supported by the observation that certain chronic infections seemed to predispose to cancer development.

\section{The Discovery and Contribution}

Fibiger's discovery began with a chance observation made in 1913\index{1913|textbf}. While examining rats from the sewers of Copenhagen, Fibiger noticed that several animals had unusual tumors in their stomachs. Upon microscopic examination, he found that the tumors contained the larvae of a nematode parasite that he identified as \emph{Spiroptera neoplastica} (later reclassified as \emph{Gongylonema neoplasticum})\index{Gongylonema neoplasticum|textbf}. This was a remarkable finding: here was a parasite that appeared to be directly causing cancer in its host.

Fibiger conducted a series of experiments to establish the causal relationship between the parasite and the stomach tumors\index{Experimental design|textbf}. He fed cockroaches (the intermediate host of the parasite) to laboratory rats, and showed that the rats developed stomach tumors at the sites where the larvae burrowed into the stomach wall. The tumors were typically squamous cell carcinomas or papillomas, and they developed at the sites of larval penetration. Fibiger demonstrated that the larvae caused chronic irritation and hyperplasia (excessive growth) of the gastric epithelium, and that this chronic irritation progressed to dysplasia (disordered growth) and ultimately to frank carcinoma\index{Hyperplasia|textbf}.

The experimental design was elegant and the results were compelling. Fibiger showed a clear dose-response relationship: rats that received more larvae developed more tumors\index{Dose-response|textbf}. He demonstrated that the tumors arose specifically at the sites of larval invasion, and that rats that were not infected with the parasite did not develop stomach tumors. He also showed that the parasite underwent a complex life cycle involving cockroaches as intermediate hosts and rats as definitive hosts.

Fibiger published his findings in a series of papers between 1913 and 1919, and his work attracted international attention\index{Publications|textbf}. The idea that a parasite could cause cancer was theoretically exciting because it suggested that cancer might be an infectious disease---a possibility that had enormous implications for prevention and treatment. The Nobel Committee, impressed by the rigor of Fibiger's experiments and the significance of his findings, awarded him the prize in 1926.

However, subsequent research cast doubt on the direct carcinogenicity of the parasite\index{Reassessment|textbf}. In the 1930s and 1940s, investigators showed that \emph{Gongylonema neoplasticum} alone was insufficient to cause cancer in rats, and that other factors---including vitamin A deficiency and concurrent infection with other parasites---were necessary for tumor development\index{Vitamin A deficiency|textbf}. The parasite was found to cause chronic irritation and hyperplasia, but it did not directly cause the genetic mutations that are now understood to be the true drivers of cancer development.

The reassessment of Fibiger's work was not a repudiation of his experimental findings but rather a refinement of his conclusions. The \emph{Gongylonema} parasite does cause chronic irritation and epithelial hyperplasia in the stomach, and these changes can, under certain conditions, progress to cancer. But the parasite is better understood as a co-factor or promoting agent rather than a direct carcinogen\index{Tumor promoter|textbf}. This distinction---between a causative agent and a promoting agent---is now recognized as fundamental to understanding carcinogenesis.

\section{Impact on Modern Medicine}

Fibiger's work, despite its limitations, had a lasting impact on cancer research and the understanding of infectious causes of cancer\index{Infectious cancer|textbf}. The principle that chronic infection and inflammation can predispose to cancer development is now well established. Chronic infections are estimated to be responsible for approximately 15--20\% of all cancers worldwide\index{Cancer epidemiology|textbf}.

The most dramatic modern confirmation of the infectious theory of cancer came with the discovery of the role of \emph{Helicobacter pylori} in gastric cancer\index{Helicobacter pylori|textbf}. Barry Marshall and Robin Warren, who received the Nobel Prize in 2005, demonstrated that \emph{H. pylori} infection causes chronic gastritis, which can progress through a well-defined sequence---atrophic gastritis, intestinal metaplasia, dysplasia, and adenocarcinoma---to gastric cancer\index{Barry Marshall|textbf}. This sequence is strikingly reminiscent of the progression from chronic irritation to carcinoma that Fibiger described a century earlier.

Human papillomavirus (HPV) infection, which causes cervical cancer, and hepatitis B and C virus infections, which cause liver cancer, are additional examples of the principle that chronic infection promotes carcinogenesis\index{HPV|textbf}. The development of vaccines against HPV and hepatitis B represents the ultimate fulfillment of the parasitic theory of cancer: if infections cause cancer, then preventing the infection should prevent the cancer. The HPV vaccine, which has the potential to virtually eliminate cervical cancer worldwide, is a notable triumph of preventive medicine in the twenty-first century.

Fibiger's work also contributed to the broader understanding of tumor promotion---the concept that cancer development is a multi-step process in which an initiating event (such as a genetic mutation) is followed by promoting events (such as chronic inflammation) that drive the progression from normal to malignant tissue\index{Tumor promotion|textbf}. The modern understanding of the tumor microenvironment---the complex ecosystem of immune cells, blood vessels, fibroblasts, and inflammatory mediators that surrounds and interacts with tumor cells---extends Fibiger's observation that the tissue environment plays a critical role in determining whether a cell becomes cancerous\index{Tumor microenvironment|textbf}.

\section{Legacy}

Johannes Fibiger's legacy is complex\index{Fibiger, Johannes|textbf}. He is sometimes cited as an example of a Nobel Prize that was awarded on the basis of findings that were later shown to be partially incorrect, but this characterization is unfair. Fibiger's experimental observations were accurate: the parasite does cause chronic irritation, hyperplasia, and, under certain conditions, cancer in the stomach. What changed was the interpretation of these findings in light of a deeper understanding of the molecular mechanisms of carcinogenesis\index{Molecular mechanisms|textbf}.

More importantly, Fibiger's work raised fundamental questions about the causes of cancer that remain relevant today\index{Cancer prevention|textbf}. The idea that external agents---whether parasites, bacteria, viruses, chemicals, or radiation---can initiate or promote the development of cancer is now the cornerstone of cancer prevention. The search for infectious causes of cancer, which Fibiger's work helped to inspire, has yielded some of the major discoveries in modern oncology.

Fibiger was a meticulous experimentalist and a careful observer, and his work exemplifies the importance of animal models in cancer research\index{Animal models|textbf}. While the specific conclusions of his work have been refined, the approach he pioneered---the use of experimental models to study the relationship between environmental exposures and cancer development---remains the foundation of cancer research. The principles of experimental pathology that Fibiger taught and practiced continue to guide investigators in the ongoing quest to understand and prevent cancer\index{Experimental pathology|textbf}. Fibiger died in Copenhagen on January 30, 1928, but his work continues to be studied and discussed as a fascinating chapter in the history of cancer research.

\section{Modern Developments}

Fibiger's work on parasite-induced cancer, though initially controversial, has been validated by modern oncology\index{Oncology|textbf}. The International Agency for Research on Cancer (IARC) classifies several parasites as carcinogenic: \emph{Schistosoma haematobium} (bladder cancer), \emph{Opisthorchis viverrini} and \emph{Clonorchis sinensis} (cholangiocarcinoma)\index{Schistosoma|textbf}. Understanding of chronic inflammation as a cancer driver has led to anti-inflammatory strategies for cancer prevention. \emph{Helicobacter pylori} (Nobel 2005) is now recognized as a Class I carcinogen for gastric cancer\index{Class I carcinogen|textbf}.
